\section{Introduzione}

Il presente elaborato si propone di condurre un'analisi comparativa delle prestazioni ottenibili parallelizzando un medesimo problema computazionale mediante i tre paradigmi affrontati nel corso: la programmazione a memoria condivisa con OpenMP, la programmazione a memoria distribuita con MPI e la programmazione su acceleratori grafici con CUDA.
L'obiettivo non è ridurre il tempo di esecuzione né competere con le implementazioni di riferimento,
bensì misurare e interpretare in modo rigoroso il comportamento di ciascun paradigma,
mettendo in relazione i risultati sperimentali con i modelli analitici del calcolo parallelo.

Come caso di studio è stato scelto il calcolo dell'\textit{\textbf{insieme di Mandelbrot}}. La scelta non è casuale.
Il valore di ciascun pixel dipende esclusivamente dalle proprie coordinate, e il problema è quindi \textit{\textbf{imbarazzantemente parallelo}} (embarrassingly parallel):
i pixel possono essere calcolati in totale indipendenza, senza che le unità di calcolo debbano scambiarsi dati.
Questa apparente semplicità nasconde tuttavia una difficoltà tutt'altro che trascurabile: il costo del singolo pixel dipende dal
numero di iterazioni necessarie a determinarne l'appartenenza all'insieme.
I punti interni raggiungono il limite massimo di iterazioni, mentre quelli esterni divergono, e quindi terminano, per lo più dopo pochi passi.
La distribuzione del carico che ne risulta è fortemente disomogenea, e non è nota prima di aver calcolato l'immagine almeno una volta.

È questa irregolarità a costituire il nucleo analitico dell'elaborato. La tesi che attraversa l'intero lavoro
è che un'unica causa, la variabilità del numero di iterazioni per pixel, si manifesti in forme diverse a seconda dell'architettura.
Sulla CPU si presenta come sbilanciamento del carico (load imbalance) fra i thread o i processi che si spartiscono l'immagine.
Sulla GPU invece si presenta come divergenza dei warp (warp divergence): i thread eseguiti insieme in uno stesso gruppo, il warp, terminano solo quando termina il più lento.
Riconoscere questa radice comune permette di trattare i tre esperimenti non come esercizi scollegati, ma come approcci alternativi a un medesimo problema.

Le prestazioni di ciascun paradigma sono misurate con metriche consolidate: speedup ed efficienza,
affiancati dalle leggi di scalabilità e da indici di sbilanciamento ricavati dal profilo di carico seriale.
La legge di Amdahl fornisce il limite allo speedup imposto dalla frazione seriale del programma, mentre la metrica di Karp–Flatt,
osservata al crescere del grado di parallelismo, aiuta a distinguere se la perdita di efficienza derivi da una componente
seriale oppure dallo sbilanciamento del carico e dalla comunicazione.

Proprio perché privo di comunicazione, tuttavia, il calcolo dell'insieme di Mandelbrot non sollecita
il regime in cui le prestazioni sono dominate dallo scambio e dalla sincronizzazione dei dati.
Un complemento naturale sarebbe il Buddhabrot, che condivide la stessa iterazione. Anziché però scrivere un solo pixel per punto,
il Buddhabrot incrementa lungo l'orbita di ciascun punto divergente numerose celle di un istogramma condiviso, introducendo contesa sulla memoria.
Non è stato affrontato in questo lavoro, e viene ripreso fra gli sviluppi (sezione~\ref{sec:conclusioni}).

Quanto ai confini del lavoro, l'elaborato non offre una trattazione della geometria frattale né dei fondamenti matematici dell'insieme: la sezione~\ref{sec:frattali} li richiama solo nella misura in cui spiegano l'origine dell'irregolarità del carico.
Sul piano sperimentale, le misure dipendono dalle risorse di calcolo disponibili, descritte nella sezione~\ref{sec:ambiente}. I valori assoluti contano quindi meno degli andamenti e del confronto fra paradigmi, e particolare cura è dedicata alla metodologia di misura e alla riproducibilità dei risultati.

\subsection{Ambiente sperimentale}
\label{sec:ambiente}

Tutte le misure sono state eseguite come job batch su un cluster HPC gestito
dallo scheduler SLURM (versione $23.11.6$), con sistema operativo Rocky
Linux~$8.9$. Si riportano di seguito le caratteristiche dell'ambiente che
incidono sulla progettazione degli esperimenti e sull'interpretazione dei
risultati.

\paragraph{Nodo di calcolo.}
L'utenza impiegata è abilitata a tre partizioni (\texttt{ulow},
\texttt{only-one-gpu} e \texttt{ice4hpc}), tutte associate al medesimo nodo,
\texttt{gnode01}; le campagne CPU sono eseguite su \texttt{ulow}, quelle GPU su
\texttt{only-one-gpu}. Il nodo monta due processori AMD EPYC~9354 (architettura
Zen~4) da $32$ core ciascuno, con estensioni AVX2 e AVX-512. Il simultaneous
multithreading è disattivato, per cui i $64$ core visibili al sistema
coincidono con i core fisici. Ogni socket costituisce un nodo NUMA con circa
$377$~GiB di memoria, per un totale di $755$~GiB; il sistema riporta per
l'accesso alla memoria dell'altro socket un costo relativo di $32$, contro $10$
per la memoria locale.

Da questa configurazione discendono tre conseguenze. Nessuna esecuzione può
estendersi su più nodi, per cui anche MPI opera su memoria condivisa
(sezione~\ref{sec:mpi}). Il tempo seriale $T_1$ e i tempi paralleli
$T_\parallel(p)$ sono misurati per costruzione sullo stesso hardware. Il nodo è
infine condiviso con i job di altri utenti: la partizione non prevede
preemption né condivisione dei core allocati, ma banda di memoria e cache di
ultimo livello restano contese, ed è per questa ragione che ogni configurazione
viene ripetuta e se ne riporta il tempo minimo.

\paragraph{Acceleratore grafico.}
Il nodo ospita otto schede NVIDIA L40S, ma la partizione dedicata alla GPU ne
assegna una sola per job: tutte le misure CUDA si riferiscono pertanto a una
singola scheda. La scheda ha architettura Ada Lovelace, compute capability
$8.9$ e $46\,068$~MiB di memoria rilevati da \texttt{nvidia-smi}; opera con il
driver $545.23.08$ ed è collegata all'host tramite un link PCIe~4.0~x16.
L'accesso ai performance counter della GPU è riservato agli amministratori, e le
misure con il profiler \texttt{ncu} restituiscono l'errore
\texttt{ERR\_NVGPUCTRPERM}: la divergenza dei warp, che richiederebbe un
profiler, è quindi stimata con un proxy deterministico
(sezione~\ref{sec:cuda-metriche}).

\paragraph{Software.}
Il compilatore è GCC $8.5.0$, fornito dal sistema; MPI è OpenMPI $4.1.6$ e CUDA
è il toolkit $12.3.2$, mantenuto invariato per l'intera campagna, poiché un
toolkit diverso produrrebbe un binario diverso e tempi GPU non confrontabili.
Ogni job compila i sorgenti sul nodo di calcolo prima di misurare: il build usa
\texttt{-march=native}, e una compilazione eseguita altrove produrrebbe codice
ottimizzato per una CPU diversa da quella su cui si misura. I rank MPI sono
lanciati con \texttt{mpirun} anziché con \texttt{srun}, a causa di
un'incompatibilità fra il client PMIx di OpenMPI $4.1.6$ e il server PMIx
esposto da SLURM (sezione~\ref{sec:mpi-codice}).
